% ============================================================
% scriptorium — ukázkový dokument (showcase všech prostředí)
% Engine: XeLaTeX.  Kompilace z této složky:
%   xelatex -interaction=nonstopmode sample.tex   (2× kvůli obsahu)
% Sdílená preambule: ../preamble/preamble.tex
% ============================================================
\documentclass[11pt, a4paper, oneside]{book}
% ============================================================
% Sdílená preambule pro český akademický / technický TeX dokument
% Engine: XeLaTeX (vyžadováno kvůli fontspec/polyglossia)
% ============================================================

% --- Jazyk a fonty ---
\usepackage{fontspec}
\usepackage{polyglossia}
\setmainlanguage{czech}
% BasicTeX nemá fonty registrované v macOS — načítáme přímo soubory přes kpathsea
\setmainfont{lmroman10-regular.otf}[
  BoldFont = lmroman10-bold.otf,
  ItalicFont = lmroman10-italic.otf,
  BoldItalicFont = lmroman10-bolditalic.otf,
  Ligatures = TeX ]
\setsansfont{lmsans10-regular.otf}[
  BoldFont = lmsans10-bold.otf,
  ItalicFont = lmsans10-oblique.otf,
  BoldItalicFont = lmsans10-boldoblique.otf,
  Ligatures = TeX ]
\setmonofont{lmmonolt10-regular.otf}[
  BoldFont = lmmonolt10-bold.otf,
  ItalicFont = lmmonolt10-oblique.otf ]

% --- Matematika ---
\usepackage{amsmath}
\usepackage{amssymb}
\usepackage{mathtools}

% --- Rozvržení stránky ---
\usepackage[a4paper, margin=2.5cm]{geometry}
\usepackage{parskip}
\usepackage{fancyhdr}
\pagestyle{fancy}
\fancyhf{}
\fancyhead[LE]{\small\leftmark}
\fancyhead[RO]{\small\rightmark}
\fancyfoot[C]{\thepage}
\renewcommand{\headrulewidth}{0.4pt}

% --- Barvy, grafika, boxy ---
\usepackage{xcolor}
\usepackage{tikz}
\usetikzlibrary{shapes.geometric, arrows.meta, positioning, calc, fit}
\usepackage[skins, breakable, theorems]{tcolorbox}
\usepackage{enumitem}
\usepackage[autostyle]{csquotes}
\usepackage{booktabs}
\usepackage{multirow}
\usepackage{tabularx}
\usepackage{longtable}
\usepackage{needspace}

% --- Barevná paleta ---
\definecolor{pojemModra}{RGB}{21, 101, 192}
\definecolor{prikladSeda}{RGB}{84, 110, 122}
\definecolor{otazkaOranzova}{RGB}{230, 126, 34}
\definecolor{reseniZelena}{RGB}{46, 125, 50}
\definecolor{pozorCervena}{RGB}{198, 40, 40}
\definecolor{kucharkaFialova}{RGB}{106, 27, 154}

% ============================================================
% Makra a prostředí (jednotný styl pro všechny kapitoly)
% ============================================================

% Anglický ekvivalent pojmu — používat při PRVNÍM výskytu každého pojmu
% Použití: úrok\en{interest}
\newcommand{\en}[1]{ \textit{(angl. #1)}}

% Seznam symbolů po vzorci
% Použití: \begin{kde} \item $FV$ — budoucí hodnota \en{future value} \end{kde}
\newenvironment{kde}%
  {\par\noindent\textit{kde:}\begin{itemize}[nosep, leftmargin=2.5em]}%
  {\end{itemize}\medskip}

% --- Pojem: definice s povinným anglickým ekvivalentem ---
% Použití: \begin{pojem}{Úrok}{interest} ... \end{pojem}
\newtcolorbox[auto counter, number within=chapter]{pojem}[2]{%
  breakable, enhanced,
  colback=pojemModra!4, colframe=pojemModra!4,
  borderline west={3pt}{0pt}{pojemModra},
  fonttitle=\bfseries, coltitle=pojemModra,
  title={Pojem~\thetcbcounter: #1 \textmd{\textit{(angl. #2)}}},
  attach title to upper={\par},
  boxrule=0pt, frame hidden, sharp corners,
  left=8pt, right=8pt, top=6pt, bottom=6pt
}

% --- Příklad: řešený příklad krok za krokem ---
% Použití: \begin{priklad}{Reálná úroková míra} ... \end{priklad}
\newtcolorbox[auto counter, number within=chapter]{priklad}[1]{%
  breakable, enhanced,
  colback=black!3, colframe=prikladSeda,
  fonttitle=\bfseries,
  title={Příklad~\thetcbcounter: #1},
  boxrule=0.8pt, arc=2pt,
  left=8pt, right=8pt, top=6pt, bottom=6pt
}

% --- Testová otázka: PŘESNÉ znění otázky ze zdroje ---
% Použití: \begin{otazka}{PF_tema_4.pdf} text otázky
%   \begin{enumerate}[label=\alph*)] \item ... \end{enumerate}
% \end{otazka}
\newtcolorbox[auto counter, number within=chapter]{otazka}[1]{%
  breakable, enhanced,
  colback=otazkaOranzova!5, colframe=otazkaOranzova,
  fonttitle=\bfseries,
  title={Testová otázka~\thetcbcounter\ \ \textmd{\small(zdroj: #1)}},
  boxrule=0.8pt, arc=2pt,
  left=8pt, right=8pt, top=6pt, bottom=6pt
}

% --- Řešení: správná odpověď + proč; proč jsou distraktory špatně ---
% Použití: \begin{reseni}{c} ... \end{reseni}   (argument = písmeno správné odpovědi)
\newtcolorbox{reseni}[1]{%
  breakable, enhanced,
  colback=reseniZelena!5, colframe=reseniZelena,
  fonttitle=\bfseries,
  title={Řešení: správná odpověď \MakeUppercase{#1})},
  boxrule=0.8pt, arc=2pt,
  left=8pt, right=8pt, top=6pt, bottom=6pt
}

% --- Odpověď: modelová odpověď na otevřenou (kontrolní) otázku ---
% Použití: \begin{odpoved} ... \end{odpoved}
\newtcolorbox{odpoved}{%
  breakable, enhanced,
  colback=reseniZelena!5, colframe=reseniZelena,
  fonttitle=\bfseries,
  title={Modelová odpověď},
  boxrule=0.8pt, arc=2pt,
  left=8pt, right=8pt, top=6pt, bottom=6pt
}

% --- Pozor: chytáky a časté chyby ---
\newtcolorbox{pozor}{%
  breakable, enhanced,
  colback=pozorCervena!5, colframe=pozorCervena,
  fonttitle=\bfseries,
  title={Pozor — častá chyba!},
  boxrule=0.8pt, arc=2pt,
  left=8pt, right=8pt, top=6pt, bottom=6pt
}

% --- Kuchařka: postup krok za krokem (čtení tabulek, integrační metody…) ---
% Použití: \begin{kucharka}{Jak číst tabulku Studentova rozdělení} \begin{enumerate} ... \end{enumerate} \end{kucharka}
\newtcolorbox[auto counter, number within=chapter]{kucharka}[1]{%
  breakable, enhanced,
  colback=kucharkaFialova!4, colframe=kucharkaFialova,
  fonttitle=\bfseries,
  title={Kuchařka~\thetcbcounter: #1},
  boxrule=0.8pt, arc=2pt,
  left=8pt, right=8pt, top=6pt, bottom=6pt
}

% --- Styly uzlů pro rozhodovací stromy (TikZ) ---
\tikzset{
  rozhodnuti/.style={diamond, aspect=2.2, draw=kucharkaFialova, fill=kucharkaFialova!8,
                     text width=9.2em, align=center, font=\small, inner sep=1pt},
  vysledek/.style={rectangle, rounded corners=3pt, draw=reseniZelena, fill=reseniZelena!8,
                   text width=10em, align=center, font=\small, inner sep=4pt},
  sipka/.style={-{Stealth}, thick},
  popisek/.style={font=\footnotesize\itshape, midway}
}

% --- Věta: matematická věta / tvrzení (rigorózní kurzy) ---
% Použití: \begin{veta}{Lagrangeova věta} ... \end{veta}
\definecolor{vetaTeal}{RGB}{0, 121, 107}
\newtcolorbox[auto counter, number within=chapter]{veta}[1]{%
  breakable, enhanced,
  colback=vetaTeal!4, colframe=vetaTeal,
  fonttitle=\bfseries, coltitle=vetaTeal,
  title={Věta~\thetcbcounter: #1},
  attach title to upper={\par},
  boxrule=0pt, frame hidden, sharp corners,
  borderline west={3pt}{0pt}{vetaTeal},
  left=8pt, right=8pt, top=6pt, bottom=6pt
}

% --- Důkaz: jednoduché prostředí s QED značkou ---
% Použití: \begin{dukaz} ... \end{dukaz}
\newenvironment{dukaz}%
  {\par\medskip\noindent\textit{Důkaz.}\ }%
  {\hfill$\square$\par\medskip}

% --- Hyperref (vždy poslední) ---
\usepackage[unicode, colorlinks=true, linkcolor=pojemModra, urlcolor=pojemModra]{hyperref}


\title{\Huge\bfseries Ukázkový kompendiový dokument\\[0.5em]
       \Large scriptorium — showcase všech prostředí\\[2em]
       \normalsize pojem \,·\, příklad \,·\, otázka \,·\, řešení \,·\, pozor \,·\, kuchařka \,·\, věta \,·\, důkaz \,·\, rozhodovací strom}
\author{}
\date{\today}

\begin{document}
\maketitle
\tableofcontents

% ============================================================
\part{Finanční ukázka}
% ============================================================

\chapter{Prostředí výkladu}
\label{ch:vyklad}

Tato kapitola ukazuje všechna prostředí, která scriptorium nabízí pro výklad
látky: definice pojmů, vzorce se seznamem symbolů, řešené příklady, varování
před chybami a~kuchařku s~postupem krok za krokem.

\section{Definice a vzorce}

Začneme základním pojmem. Anglický ekvivalent připojujeme makrem
\verb|\en{}| při prvním výskytu termínu.

\begin{pojem}{Úrok}{interest}
Úrok je cena za zapůjčení peněz vyjádřená v~penězích. Věřitel\en{creditor}
půjčí dlužníkovi\en{debtor} jistinu a~po čase dostane zpět jistinu navýšenou
o~úrok.
\end{pojem}

Budoucí hodnotu vkladu při ročním složeném úročení spočítáme takto:
\[ FV = PV \cdot (1 + r)^t \]
\begin{kde}
  \item $FV$ — budoucí hodnota\en{future value},
  \item $PV$ — současná hodnota\en{present value},
  \item $r$ — roční úroková míra (desetinné číslo),
  \item $t$ — počet let.
\end{kde}

\begin{priklad}{Budoucí hodnota vkladu}
Vložíme $PV = 10\,000$ Kč na účet s~úrokem $r = 5\,\%$ ročně na $t = 3$ roky.
\begin{enumerate}
  \item Dosadíme do vzorce: $FV = 10\,000 \cdot (1 + 0{,}05)^3$.
  \item Spočítáme závorku: $1 + 0{,}05 = 1{,}05$.
  \item Umocníme: $1{,}05^3 = 1{,}157625$.
  \item Vynásobíme: $10\,000 \cdot 1{,}157625 = 11\,576{,}25$.
\end{enumerate}
\textbf{Výsledek: 11\,576,25 Kč.}
\end{priklad}

\begin{pozor}
Úrokovou míru $5\,\%$ dosazujeme do vzorce jako desetinné číslo $0{,}05$,
nikoli jako $5$. Záměna je nejčastější chyba u~úrokových výpočtů.
\end{pozor}

\begin{kucharka}{Jak spočítat budoucí hodnotu jednorázového vkladu}
\begin{enumerate}
  \item Zapiš si jistinu $PV$, úrokovou míru $r$ a~dobu $t$.
  \item Úrokovou míru převeď na desetinné číslo ($5\,\% \to 0{,}05$).
  \item Dosaď do $FV = PV \cdot (1+r)^t$ a~vyčísli.
\end{enumerate}
\end{kucharka}

\section{Testové otázky ze zdrojových materiálů}

\begin{otazka}{ukazka\_finance.pdf}
Kolik činí budoucí hodnota vkladu $1\,000$ Kč při úroku $10\,\%$ ročně po
dvou letech (složené úročení)?
\begin{enumerate}[label=\alph*)]
  \item $1\,100$ Kč
  \item $1\,200$ Kč
  \item $1\,210$ Kč
  \item $1\,000$ Kč
\end{enumerate}
\end{otazka}
\begin{reseni}{c}
Platí $FV = 1\,000 \cdot 1{,}1^2 = 1\,000 \cdot 1{,}21 = 1\,210$ Kč.
Odpověď a) počítá jen jeden rok, b) sčítá úrok lineárně bez složeného
úročení, d) zapomíná úrok připočíst.
\end{reseni}

\begin{otazka}{ukazka\_kontrolni.pdf}
Vysvětlete vlastními slovy rozdíl mezi jednoduchým a~složeným úročením.
\end{otazka}
\begin{odpoved}
U~jednoduchého úročení se úrok počítá stále z~původní jistiny, takže roste
lineárně. U~složeného úročení se úrok připisuje k~jistině a~v~dalším období
se už úročí i~dříve připsaný úrok — jistina proto roste exponenciálně.
\end{odpoved}

% ============================================================
\part{Matematická ukázka}
% ============================================================

\chapter{Věty, důkazy a rozhodovací stromy}
\label{ch:veta}

Pro rigoróznější kapitoly nabízí scriptorium prostředí pro matematické věty,
důkaz se značkou QED a~styly uzlů pro rozhodovací stromy v~TikZ.

\begin{veta}{Součet prvních $n$ přirozených čísel}
Pro každé přirozené číslo $n$ platí
\[ 1 + 2 + \dots + n = \frac{n(n+1)}{2}. \]
\end{veta}

\begin{dukaz}
Označme součet $S = 1 + 2 + \dots + n$. Sečteme-li tutéž sumu napsanou
pozpátku, dostaneme $2S = (n+1) + (n+1) + \dots + (n+1) = n(n+1)$, odkud
$S = \frac{n(n+1)}{2}$.
\end{dukaz}

\section{Rozhodovací strom}

Tam, kde je třeba zvolit metodu, používáme rozhodovací strom se styly uzlů
\verb|rozhodnuti|, \verb|vysledek|, \verb|sipka| a~\verb|popisek|.

\begin{center}
\textbf{Jak derivovat výraz?}\\[0.8em]
\begin{tikzpicture}[node distance=1.1cm and 1.8cm]
  \node[rozhodnuti] (d1) {Je to elementární funkce z~tabulky?};
  \node[vysledek, right=of d1] (r1) {použij \textbf{tabulku derivací}, např. $(x^n)' = n\,x^{n-1}$};
  \node[rozhodnuti, below=of d1] (d2) {Je to součin $u \cdot v$?};
  \node[vysledek, right=of d2] (r2) {\textbf{pravidlo součinu}: $(uv)' = u'v + uv'$};
  \node[rozhodnuti, below=of d2] (d3) {Je to složená funkce $f(g(x))$?};
  \node[vysledek, right=of d3] (r3) {\textbf{řetízkové pravidlo}: $f'(g)\cdot g'$};
  \draw[sipka] (d1) -- node[popisek, above] {ano} (r1);
  \draw[sipka] (d1) -- node[popisek, right] {ne} (d2);
  \draw[sipka] (d2) -- node[popisek, above] {ano} (r2);
  \draw[sipka] (d2) -- node[popisek, right] {ne} (d3);
  \draw[sipka] (d3) -- node[popisek, above] {ano} (r3);
\end{tikzpicture}
\end{center}

\end{document}
